% Part: computability
% Chapter: computability-theory
% Section: equiv-ce-defs

\documentclass[../../../include/open-logic-section]{subfiles}

\begin{document}

\olfileid{cmp}{thy}{eqc}

\olsection[C. E. সেটের সংজ্ঞা]{গণনসাধ্যভাবে তালিকায়নযোগ্য সেটের
  সমতুল্য সংজ্ঞা}


গণনসাধ্যভাবে তালিকায়নযোগ্য হওয়ার অর্থ কী, নিচে তার কয়েকটি গুরুত্বপূর্ণ
সমতুল্য বক্তব্য দেওয়া হলো।

\begin{thm}
\ollabel{thm:ce-equiv}
ধরা যাক, $S$ স্বাভাবিক সংখ্যার একটি সেট। তাহলে নিচের বক্তব্যগুলি সমতুল্য:
\begin{enumerate}
\item\ollabel{case:ce} $S$ গণনসাধ্যভাবে তালিকায়নযোগ্য।
\item\ollabel{case:ran-pc} $S$ কোনো \emph{আংশিক} গণনসাধ্য অপেক্ষকের মানপরিসর।
\item\ollabel{case:ran-prim} $S$ শূন্য অথবা কোনো আদিম পুনরাবৃত্ত অপেক্ষকের মানপরিসর।
\item\ollabel{case:ce-domain} $S$ কোনো আংশিক গণনসাধ্য অপেক্ষকের
\emph{সংজ্ঞাক্ষেত্র}।
\end{enumerate}
\end{thm}

\begin{explain}
প্রথম তিনটি দফা বলে, যেকোনো অশূন্য গণনসাধ্যভাবে তালিকায়নযোগ্য সেটকে
গণনসাধ্য অপেক্ষক, আংশিক গণনসাধ্য অপেক্ষক বা আদিম পুনরাবৃত্ত অপেক্ষক—এই
তিনটির যেকোনো একটি দিয়ে সমতুল্যভাবে তালিকায়িত করা যায়। চতুর্থ দফাটি
বলে, $S$ গণনসাধ্যভাবে তালিকায়নযোগ্য হলে কোনো সূচক~$e$-এর জন্য
\[
S = \Setabs{x}{\cfind{e}(x) \fdefined}.
\]
অন্যভাবে বললে, $S$ হলো সেই সব নিবেশের সেট, যেগুলিতে $\cfind{e}$-এর গণনা
থামে। এই কারণে গণনসাধ্যভাবে তালিকায়নযোগ্য সেটকে কখনও কখনও
\emph{অর্ধ-নির্ণেয়} বলা হয়: কোনো সংখ্যা সেটে থাকলে শেষে “হ্যাঁ” পাওয়া
যায়, কিন্তু না থাকলে কখনও “না” পাওয়া যায় না!{}
\end{explain}

\begin{proof}
প্রত্যেক আদিম পুনরাবৃত্ত অপেক্ষক গণনসাধ্য এবং প্রত্যেক গণনসাধ্য অপেক্ষক
আংশিক গণনসাধ্য বলে \olref{case:ran-prim} থেকে \olref{case:ce}, এবং
\olref{case:ce} থেকে~\olref{case:ran-pc} অনুসৃত হয়। ($S$ শূন্য হলে
$S$~সর্বত্র অসংজ্ঞায়িত আংশিক গণনসাধ্য অপেক্ষকটির মানপরিসর।) আমরা যদি
দেখাই যে \olref{case:ran-pc} থেকে \olref{case:ran-prim} অনুসৃত হয়, তবে
প্রথম তিনটি দফার সমতুল্যতা প্রমাণিত হবে।

ধরা যাক, $S$ আংশিক গণনসাধ্য অপেক্ষক $\cfind{e}$-এর মানপরিসর। $S$ শূন্য
হলে প্রমাণ শেষ। অন্যথায়, $a$-কে~$S$-এর যেকোনো একটি উপাদান ধরি। ক্লিনির
স্বাভাবিক রূপের উপপাদ্য অনুসারে লেখা যায়
\[
\cfind{e}(x) \simeq U(\umin{s}{T(e, x, s)}).
\]
% BN-SRC-147 TARGET-CORRECTION-BEGIN
\par\smallskip\noindent\textnormal{\small\textbf{উৎস-সংশোধন (BN-SRC-147):} হিমায়িত ইংরেজি প্রদর্শনে আংশিক অপেক্ষকের স্বাভাবিক রূপে “$=$” আছে, কিন্তু উদ্ধৃত স্বাভাবিক-রূপ উপপাদ্যটি যুগপৎ সংজ্ঞায়িততার সমতা “$\simeq$” দেয়। বাংলা প্রদর্শনে সেই সঠিক চিহ্ন রাখা হয়েছে।} % BN-SRC-147 TARGET-CORRECTION-END
বিশেষত, $\cfind{e}(x) \fdefined$ এবং $= y$ যদি এবং কেবল যদি এমন একটি
$s$ থাকে, যাতে $T(e, x, s)$ এবং $U(s) = y$। $f(z)$-কে সংজ্ঞায়িত করি:
\[
f(z) = \begin{cases}
  U((z)_1) & \text{যদি $T(e, (z)_0, (z)_1)$} \\
  a        & \text{অন্যথায়।}
\end{cases}
\]
তাহলে $f$ আদিম পুনরাবৃত্ত, কারণ $T$ ও $U$ আদিম পুনরাবৃত্ত।
\iftag{TMs}{টুরিং যন্ত্রের ভাষায়, $z$ যদি এমন একটি জোড়া
  $\tuple{(z)_0, (z)_1}$-কে সংকেতায়িত করে যেখানে $(z)_1$ হলো নিবেশ
  $(z)_0$-এ যন্ত্র~$M_e$-এর থামা গণনা, তবে $f$ গণনাটির নির্গম ফেরত দেয়;
  অন্যথায়~$a$ ফেরত দেয়।}

আমাদের দেখাতে হবে যে $S$ হলো~$f$-এর মানপরিসর, অর্থাৎ যেকোনো স্বাভাবিক
সংখ্যা~$y$-এর জন্য $y \in S$ যদি এবং কেবল যদি সেটি~$f$-এর মানপরিসরে
থাকে। সম্মুখ দিকে, ধরা যাক $y \in S$। তাহলে $y$, $\cfind{e}$-এর
মানপরিসরে আছে; তাই কোনো $x$ ও~$s$-এর জন্য $T(e,x,s)$ সত্য এবং
$U(s) = y$। কিন্তু তখন $y = f(\tuple{x,s})$। বিপরীত দিকে, ধরা যাক $y$
হলো~$f$-এর মানপরিসরের অন্তর্ভুক্ত। তাহলে হয় $y = a$, নয়তো কোনো~$z$-এর
জন্য $T(e,(z)_0,(z)_1)$ এবং $U((z)_1) = y$। পরের ক্ষেত্রে
$\cfind{e}((z)_0) \fdefined = y$ বলে উভয় ক্ষেত্রেই~$y$, $S$-তে আছে।
% BN-SRC-148 TARGET-CORRECTION-BEGIN
\par\smallskip\noindent\textnormal{\small\textbf{উৎস-সংশোধন (BN-SRC-148):} হিমায়িত ইংরেজি বিপরীত দিকের যুক্তিতে $\cfind{e}(x)$ লেখা হয়েছে, যদিও এই দিকটিতে সাক্ষী~$z$ এবং গণনার নিবেশ~$(z)_0$; আগের সম্মুখ দিকের~$x$ এখানে আবদ্ধ নয়। বাংলা পাঠে $\cfind{e}((z)_0)$ রাখা হয়েছে।} % BN-SRC-148 TARGET-CORRECTION-END

($\cfind{e}(x) \fdefined = y$ সংকেতটির অর্থ, “$\cfind{e}(x)$ সংজ্ঞায়িত
এবং~$y$-এর সমান।” একই কাজে $\cfind{e}(x) = y$ ব্যবহার করা যেত, কিন্তু
অতিরিক্ত তিরচিহ্নটি মাঝে মাঝে মনে করিয়ে দেয় যে আমরা একটি আংশিক অপেক্ষক
নিয়ে কাজ করছি।)

\olref{thm:ce-equiv}-এর প্রমাণ শেষ করতে \olref{case:ce} ও
\olref{case:ce-domain}-এর সমতুল্যতা দেখালেই যথেষ্ট। প্রথমে দেখাই,
\olref{case:ce} থেকে~\olref{case:ce-domain} অনুসৃত হয়। সেটটি শূন্য হলে
সর্বত্র অসংজ্ঞায়িত আংশিক গণনসাধ্য অপেক্ষকের সংজ্ঞাক্ষেত্রই সেটটি; তাই
এবার অশূন্য ক্ষেত্রটি ধরি।
% BN-SRC-149 TARGET-CORRECTION-BEGIN
\par\smallskip\noindent\textnormal{\small\textbf{উৎস-সংযোজন (BN-SRC-149):} গণনসাধ্যভাবে তালিকায়নযোগ্য সেটের সংজ্ঞায় শূন্য সেটটি আলাদাভাবে অন্তর্ভুক্ত, কিন্তু হিমায়িত ইংরেজি প্রমাণটি সঙ্গে সঙ্গে সেটটিকে কোনো সর্বত্র-সংজ্ঞায়িত অপেক্ষকের মানপরিসর ধরে—যা শূন্য সেটের জন্য অসম্ভব। বাংলা পাঠে সর্বত্র অসংজ্ঞায়িত আংশিক গণনসাধ্য অপেক্ষক দিয়ে শূন্য ক্ষেত্রটি আগে নিষ্পন্ন করা হয়েছে।} % BN-SRC-149 TARGET-CORRECTION-END
ধরা যাক, $S$ কোনো গণনসাধ্য অপেক্ষক~$f$-এর মানপরিসর, অর্থাৎ
\[
S = \Setabs{y}{\text{কোনো $x$-এর জন্য } f(x) = y}.
\]
ধরি,
\[
g(y) = \umin{x}{(f(x) = y)}.
\]
তাহলে $g$ একটি আংশিক গণনসাধ্য অপেক্ষক, এবং কোনো~$x$-এর জন্য
$f(x) = y$ হলে এবং কেবল হলেই $g(y)$ সংজ্ঞায়িত। অন্যভাবে বললে, $g$-এর
সংজ্ঞাক্ষেত্র হলো~$f$-এর মানপরিসর। \iftag{TMs}{টুরিং যন্ত্রের ভাষায়:
কোনো টুরিং যন্ত্র~$F$, $S$-এর উপাদানগুলি তালিকায়িত করলে $G$-কে এমন
টুরিং যন্ত্র ধরি, যা প্রদত্ত উপাদানটি সেটে আছে কি না দেখতে~$F$-এর
নির্গমগুলির মধ্যে অনুসন্ধান করে~$S$-কে অর্ধ-নির্ণয় করে; উপাদানটি থাকলে
থামে, না থাকলে চিরকাল খুঁজতে থাকে।}{}

শেষে, \olref{case:ce-domain} থেকে~\olref{case:ce} দেখাতে ধরা যাক,
$S$~আংশিক গণনসাধ্য অপেক্ষক~$\cfind{e}$-এর সংজ্ঞাক্ষেত্র, অর্থাৎ
\[
S = \Setabs{x}{\cfind{e}(x) \fdefined}.
\]
$S$ শূন্য হলে প্রমাণ শেষ; অন্যথায়, $a$-কে~$S$-এর যেকোনো উপাদান ধরি।
$f$-কে সংজ্ঞায়িত করি:
\[
f(z) = \begin{cases}
(z)_0 & \text{যদি $T(e,(z)_0,(z)_1)$} \\
a & \text{অন্যথায়।}
\end{cases}
\]
তাহলে উপরের মতোই কোনো সংখ্যা $x$, $f$-এর মানপরিসরে আছে যদি এবং কেবল
যদি $\cfind{e}(x) \fdefined$, অর্থাৎ যদি এবং কেবল যদি $x \in S$।
\iftag{TMs}{টুরিং যন্ত্রের ভাষায়: কোনো যন্ত্র $M_e$, $S$-কে অর্ধ-নির্ণয়
করলে সম্ভাব্য সব টুরিং-যন্ত্র গণনার মধ্য দিয়ে গিয়ে এবং থামা গণনাগুলির
সংশ্লিষ্ট নিবেশগুলি ফেরত দিয়ে~$S$-এর উপাদানগুলি তালিকায়িত করা যায়।}{}
\end{proof}

\olref{thm:ce-equiv}-এর দফা~\olref{case:ce-domain} গণনসাধ্যভাবে
তালিকায়নযোগ্য সেটগুলিকে তালিকায়িত করার সুবিধাজনক উপায় দেয়: প্রত্যেক~$e$-এর
জন্য $W_e$ দিয়ে $\cfind{e}$-এর সংজ্ঞাক্ষেত্র বোঝাই, অর্থাৎ
\[
W_e = \Setabs{x}{\cfind{e}(x) \fdefined}.
\]
তাহলে $A$ যেকোনো গণনসাধ্যভাবে তালিকায়নযোগ্য সেট হলে কোনো~$e$-এর জন্য
$A = W_e$।

নিচে গণনসাধ্যভাবে তালিকায়নযোগ্য সেটের আরও একটি চরিত্রায়ন দেওয়া হলো।

\begin{thm}
\ollabel{thm:exists-char}
কোনো সেট $S$ গণনসাধ্যভাবে তালিকায়নযোগ্য যদি এবং কেবল যদি এমন একটি
গণনসাধ্য সম্বন্ধ $R(x,y)$ থাকে, যাতে
\[
S = \Setabs{ x }{ \lexists[y][R(x,y)] }.
\]
\end{thm}

\begin{proof}
সম্মুখ দিকে, ধরা যাক $S$ গণনসাধ্যভাবে তালিকায়নযোগ্য। তাহলে কোনো $e$-এর
জন্য $S = W_e$। এই~$e$-এর জন্য~$S$-কে লেখা যায়
\[
S = \Setabs{ x }{ \lexists[y][T(e, x, y)] }.
\]
বিপরীত দিকে, ধরা যাক $S = \Setabs{ x }{
  \lexists[y][R(x, y)] }$। $f$-কে সংজ্ঞায়িত করি:
\[
f(x) \simeq \umin{y}{R(x, y)}.
\]
তাহলে $f$ আংশিক গণনসাধ্য এবং~$S$ হলো~$f$-এর সংজ্ঞাক্ষেত্র।
\end{proof}


\end{document}
